\documentclass[12pt]{article}
\usepackage{amsmath, amssymb, graphicx, hyperref, geometry}
\geometry{a4paper, margin=1in}

\title{\textbf{GANs from Scratch: Part 1 -- The Math of Adversarial Training}}
\author{Build in Public Series}
\date{\today}

\begin{document}
\maketitle

\begin{center}
    \includegraphics[width=\textwidth]{Post-1-Banner.png}
\end{center}

\section{Introduction}
In 2014, Ian Goodfellow introduced a radical idea: instead of training a single network to generate data, train \textit{two} networks that compete against each other. A Generator that tries to produce fake data, and a Discriminator that tries to distinguish real from fake. This adversarial game, formalized as a minimax optimization, launched the field of Generative Adversarial Networks (GANs) and fundamentally reshaped generative modeling.

In Part 1 of this mini-series, we will deconstruct the mathematical core of GANs: the minimax objective, Jensen-Shannon divergence, the proof of the optimal discriminator, and why training GANs is equivalent to searching for a Nash equilibrium.

\section{The Minimax Game}
The GAN framework consists of two players:
\begin{itemize}
    \item \textbf{Generator $G$}: Takes a random noise vector $z \sim p_z(z)$ and maps it to the data space. Its goal is to produce samples that are indistinguishable from real data.
    \item \textbf{Discriminator $D$}: Takes a data sample $x$ and outputs a probability $D(x) \in [0, 1]$ that $x$ came from the real data distribution rather than from $G$.
\end{itemize}

The two networks play a minimax game with the value function:
\begin{equation}
    \min_G \max_D \; V(D, G) = \mathbb{E}_{x \sim p_{\text{data}}(x)}[\log D(x)] + \mathbb{E}_{z \sim p_z(z)}[\log(1 - D(G(z)))]
\end{equation}

The Discriminator wants to \textit{maximize} this objective: assign high probability to real data ($D(x) \to 1$) and low probability to fake data ($D(G(z)) \to 0$). The Generator wants to \textit{minimize} it: make $D(G(z)) \to 1$ so that $\log(1 - D(G(z))) \to -\infty$.

\section{Jensen-Shannon Divergence}
Why does this particular objective work? It turns out that the minimax game implicitly minimizes the \textit{Jensen-Shannon divergence} between the real data distribution $p_{\text{data}}$ and the generated distribution $p_g$.

The JSD is defined as:
\begin{equation}
    \text{JSD}(p_{\text{data}} \| p_g) = \frac{1}{2} D_{\text{KL}}\left(p_{\text{data}} \;\middle\|\; \frac{p_{\text{data}} + p_g}{2}\right) + \frac{1}{2} D_{\text{KL}}\left(p_g \;\middle\|\; \frac{p_{\text{data}} + p_g}{2}\right)
\end{equation}

Unlike the KL divergence, JSD is symmetric and bounded: $0 \leq \text{JSD} \leq \log 2$. When $p_g = p_{\text{data}}$, $\text{JSD} = 0$, and the Generator has perfectly learned the data distribution.

\section{The Optimal Discriminator}
For a fixed Generator $G$, the optimal Discriminator $D^*$ can be derived analytically. Taking the functional derivative of $V(D, G)$ with respect to $D(x)$ and setting it to zero:

\begin{equation}
    D^*(x) = \frac{p_{\text{data}}(x)}{p_{\text{data}}(x) + p_g(x)}
\end{equation}

This is an elegant result. At optimality, the Discriminator outputs the \textit{density ratio} between the real and generated distributions. When $p_g = p_{\text{data}}$, we get $D^*(x) = \frac{1}{2}$ everywhere --- the Discriminator is reduced to random guessing.

Substituting $D^*$ back into the value function yields:
\begin{equation}
    V(D^*, G) = -\log 4 + 2 \cdot \text{JSD}(p_{\text{data}} \| p_g)
\end{equation}

This proves that training the Generator against the optimal Discriminator is equivalent to minimizing the JSD between $p_{\text{data}}$ and $p_g$.

\section{Training Dynamics and Nash Equilibrium}
GAN training alternates between updating $D$ and updating $G$. In game theory, the solution to a minimax game is a \textit{Nash equilibrium} --- a state where neither player can improve by unilaterally changing their strategy.

For GANs, the Nash equilibrium occurs when $p_g = p_{\text{data}}$ and $D(x) = \frac{1}{2}$ for all $x$. However, reaching this equilibrium is notoriously difficult in practice:

\begin{itemize}
    \item \textbf{Mode collapse}: The Generator learns to produce only a few high-confidence samples rather than covering the full data distribution.
    \item \textbf{Training instability}: If the Discriminator becomes too strong too quickly, the Generator receives vanishing gradients ($\log(1 - D(G(z))) \to 0$) and cannot learn.
    \item \textbf{Oscillation}: The two networks may chase each other around the equilibrium without converging.
\end{itemize}

In practice, Goodfellow proposed a modification: instead of minimizing $\log(1 - D(G(z)))$, the Generator maximizes $\log(D(G(z)))$. This provides stronger gradients early in training when $D(G(z)) \approx 0$.

\section{Up Next}
With the mathematical foundations laid, we understand \textit{what} a GAN is optimizing and \textit{why} it works. In Part 2, we will implement these ideas in pure PyTorch --- building both a Vanilla GAN and a Deep Convolutional GAN (DCGAN) from scratch, with careful attention to the architectural choices that make training stable.

\end{document}
